% EMMET paper - LaTeX source
% Compile: pdflatex emmet_paper.tex (twice for refs)
\documentclass[11pt,a4paper]{article}

\usepackage[utf8]{inputenc}
\usepackage[T1]{fontenc}
\usepackage{lmodern}
\usepackage{microtype}
\usepackage[margin=2.5cm]{geometry}
\usepackage{amsmath,amssymb}
\usepackage{graphicx}
\usepackage{booktabs}
\usepackage{array}
\usepackage{hyperref}
\usepackage{xcolor}
\usepackage{enumitem}
\usepackage{caption}

\hypersetup{
  colorlinks=true,
  linkcolor=blue!50!black,
  citecolor=blue!50!black,
  urlcolor=blue!50!black,
  pdftitle={EMMET: A Potential-Field Routing Heuristic with Adaptive Congestion Thermostat and Loss Memory},
  pdfauthor={Carlos Lopez}
}

\title{EMMET: A Potential-Field Routing Heuristic\\
with Adaptive Congestion Thermostat and Loss Memory}

\author{Carlos L\'opez\\
\small Independent Researcher\\
\small \texttt{carlos@codix.cat}\\
\small \href{https://github.com/Carloscodix/EMMET}{\texttt{github.com/Carloscodix/EMMET}}}

\date{\today}

\begin{document}

\maketitle

\begin{abstract}
We present EMMET, an adaptive routing algorithm grounded in classical
mechanics and thermodynamics. Network packets are modeled as physical
particles navigating a composite potential field with three terms:
distance to destination, instantaneous link congestion, and decayed
historical packet loss. The TTL is interpreted as physical energy;
packets that exhaust their budget are considered to undergo thermal
dissipation. The loss term carries a frozen warm-up snapshot subject to
half-life decay, modeling persistent thermal memory of past failures.
An $\varepsilon$-greedy stochastic term implements local exploration. A
global congestion thermostat adapts the congestion weight $\beta$
according to the mean normalized network load.

We characterize a phase transition in the algorithm's behavior as a
function of network density: below a critical threshold
$\rho_c \in [0.15, 0.30]$ on Erd\H{o}s--R\'enyi graphs, the potential
field collapses into topological dead ends. Above the threshold, EMMET
achieves near-zero packet loss in the dense regime (mean below 0.2
packets per 200-step run for $\rho \geq 0.30$ on $n=20$). We further
identify a $\beta$ sweet spot at $\beta = 3.5$--$4.0$ followed by a
field saturation regime, and the empirical relationship between mechanisms:
thermal memory and adaptive $\beta$ are \textbf{synergistic}, while
local lookahead $h{=}2$ and thermal memory are \textbf{substitutes}.

EMMET is validated against three baselines (Shortest Path, ECMP, LASP)
on synthetic Erd\H{o}s--R\'enyi graphs ($N \in \{20, 50, 100\}$) and
real Internet topologies (Abilene, GEANT) from the Internet Topology
Zoo, across three batteries totaling 28{,}800 simulations
(100 seeds per scenario for $n=20$, $n=50$, and real topologies; 50
seeds for $n=100$ due to compute cost). The full system (adaptive
$\beta$ + thermal) reduces packet loss by 55\,\% to 65\,\% versus LASP
on synthetic networks and 54.5\,\% on GEANT, with a 12.3\,\%
reduction on the small Abilene topology that previous variants could
not improve. Three independent adversarial code audits identified and
corrected methodological errors.
\end{abstract}

\textbf{Keywords:} adaptive routing, potential field, phase transition,
thermal dynamics, complex networks, network science.

\section{Introduction}

The standard approach to network routing --- minimize path length under
some cost metric --- treats the network as a static optimization
problem. The packet follows the shortest path and arrives, or it does
not.

Real networks are dynamic systems. Links congest, degrade, and recover.
A packet that blindly follows the shortest path may arrive at a
saturated link and be dropped, while an alternative longer path
delivers reliably.

This paper asks: \textit{what if a packet behaved like a physical
particle?} In classical mechanics, a particle moving through a field
does not follow the straight-line path --- it follows the path of
least potential energy, trading distance for the avoidance of
high-energy regions. Friction dissipates motion. Heat marks regions of
maximum energy exchange. A particle with finite energy that cannot
reach its destination dies of thermal dissipation.

This framing produces a concrete algorithm with measurable physical
analogues. We name it EMMET (Energy-Minimizing Multi-dimensional
Edge-based Traversal). Each routing decision minimizes a composite
potential combining distance, congestion, and loss memory, under a
globally adapted congestion weight, with stochastic exploration to
break greedy rigidity.

Our contributions are:
\begin{enumerate}[itemsep=2pt]
  \item A potential-field routing algorithm with four physically
        motivated mechanisms: warm-up loss snapshot, half-life decay,
        $\varepsilon$-greedy exploration, and adaptive $\beta$
        thermostat.
  \item Empirical characterization of a critical density threshold
        below which the potential field collapses (phase transition).
  \item Identification of a $\beta$ sweet spot followed by a field
        saturation regime.
  \item Empirical demonstration that thermal memory and adaptive
        $\beta$ are synergistic, while lookahead $h{=}2$ substitutes
        for thermal memory.
  \item Validation across 28{,}800 simulations on synthetic and real
        Internet topologies, after three adversarial code audits.
\end{enumerate}

\section{Related Work}

\textbf{Potential-field routing} has been studied for wireless sensor
and ad-hoc networks. Lenders et al.~\cite{lenders2008} propose
density-based anycast with explicit analysis of local maxima, deriving
theoretical bounds for local-minima-free operation under a single-term
potential. Their analysis does not characterize a critical density for
composite multi-term potentials.

\textbf{Geographic routing} literature~\cite{karp2000gpsr} documents
local minima in greedy forwarding and addresses them through recovery
mechanisms (face routing, perimeter routing) rather than characterizing
the failure regime.

\textbf{Backpressure routing}~\cite{tassiulas1992} routes on queue
differentials and is throughput-optimal but degrades latency.

\textbf{ECMP}~\cite{rfc2992} distributes flows hash-based among
equal-cost paths and is the industry standard for load balancing in
IP networks. \textbf{LASP} (Load-Aware Shortest Path) extends Dijkstra
with congestion-aware weights and is a standard traffic-engineering
baseline.

None of these works characterize a critical density threshold for a
composite three-term potential function combining distance, congestion,
and decayed loss memory; nor do they study mechanism interactions
(synergy vs.\ substitution) between exploration, lookahead, and
adaptive parameters. This is the gap EMMET addresses.

\section{The EMMET Model}

\subsection{Network Model}

The network is an undirected graph $G = (V, E)$ where each edge
$(u,v)$ carries:
\begin{itemize}[itemsep=1pt]
  \item $\mathrm{latency}(u,v)$ --- static propagation delay
  \item $\mathrm{capacity}(u,v)$ --- maximum simultaneous load before
        a packet is dropped
  \item $\mathrm{load}(u,v)$ --- current traffic level
  \item $\mathrm{loss}(u,v)$ --- accumulated packet loss count
\end{itemize}

\paragraph{Loss and load dynamics.} At each simulation step a single
packet is injected with a random source--destination pair. The packet
traverses the chosen path edge by edge. Before each edge transit
$\mathrm{load}(u,v) \leftarrow \mathrm{load}(u,v) + 1$; if the new
load exceeds $\mathrm{capacity}(u,v)$, the packet is dropped at that
edge, $\mathrm{loss}(u,v)$ is incremented by one, and no further hops
are attempted. There are no queues or retransmissions. After each
step, $\mathrm{load}$ on every edge decays by a factor of $0.9$,
modelling traffic that is bursty rather than permanently saturating.
This idealisation isolates the effect of the routing decision on
loss, at the cost of not modelling sustained flows; we discuss this
limitation in \S\ref{sec:limits}.

\subsection{Potential Function}

For a packet at node $u$ considering neighbor $v$ toward destination
$\mathrm{dst}$:
\begin{equation}
P(u, v, \mathrm{dst}) = \alpha \cdot d(v, \mathrm{dst})
                      + \beta_{\mathrm{eff}} \cdot
                        \frac{\mathrm{load}(u,v)}{\mathrm{capacity}(u,v)}
                      + \gamma \cdot \mathrm{loss\_snapshot}(u,v)
\end{equation}
where $\alpha$ and $\gamma$ are fixed weights, and
$\beta_{\mathrm{eff}}$ is governed by the adaptive thermostat
(\S\ref{sec:thermostat}).

\subsection{Routing Rule}

At each hop, the packet moves to the neighbor minimizing $P$,
excluding previously visited nodes. With probability $\varepsilon$, the
second-best neighbor is chosen instead --- this implements stochastic
exploration (\S\ref{sec:eps}).

\subsection{TTL as Physical Energy}

A packet is born with energy budget
$E = \mathrm{TTL\_FACTOR} \cdot |V|$ hops. Two distinct termination
conditions exist:
\begin{itemize}[itemsep=1pt]
  \item \textbf{dead\_end}: no unvisited neighbors --- topological
        local minimum
  \item \textbf{ttl\_expired}: energy exhausted --- thermal dissipation
\end{itemize}
These are tracked separately as failure modes signaling field collapse
versus insufficient energy for the topology.

\subsection{Thermal Memory: Snapshot, Decay, and Warm-up}

The $\mathrm{loss\_snapshot}$ term carries persistent information
about past failures. It is constructed in three steps:

\textbf{Warm-up phase.} Before measurement, EMMET routes
$\mathrm{warmup\_steps} = \max(20, |V| \cdot 5)$ packets and freezes
the resulting per-edge loss values into a snapshot dictionary.

\textbf{Read-only during measurement.} The snapshot is consulted at
every routing decision but is not updated by losses occurring during
measurement. This eliminates information asymmetry against the
baselines.

\textbf{Half-life decay.} After each step,
$\mathrm{snapshot}(u,v) \leftarrow \mathrm{decay} \cdot
\mathrm{snapshot}(u,v)$, with
$\mathrm{decay} = \exp(-\ln 2 / \mathrm{HALF\_LIFE})$. For
$\mathrm{HALF\_LIFE} = 100$ steps, $\mathrm{decay} \approx 0.9931$ ---
the snapshot loses 50\,\% of its value every 100 steps. This models
heat dissipation.

\subsection{Adaptive \texorpdfstring{$\beta$}{beta} Thermostat}
\label{sec:thermostat}

A fixed $\beta$ assumes uniform congestion across the network. Real
topologies have heterogeneously stressed regions. The adaptive
thermostat scales $\beta$ with the global thermal state:
\begin{equation}
\beta_{\mathrm{eff}} = \beta_{\mathrm{base}} \cdot
  \left(1 + \theta \cdot
  \left\langle \frac{\mathrm{load}(u,v)}{\mathrm{capacity}(u,v)}
  \right\rangle \right)
\end{equation}
where $\langle \cdot \rangle$ denotes the mean over all edges and
$\theta \geq 0$ is the sensitivity. With $\theta = 0$, $\beta$ is
fixed; with $\theta > 0$, $\beta$ scales up under network stress. The
thermostat introduces no learning state --- it is a stateless
feedback mechanism evaluated at routing time.

\subsection{\texorpdfstring{$\varepsilon$}{Epsilon}-greedy Stochastic Exploration}
\label{sec:eps}

With probability $\varepsilon$, the routing decision selects the
second-best neighbor by potential ranking instead of the best. This
breaks greedy rigidity, enables exploration of locally suboptimal
paths, and provides robustness to small modeling errors. We use
$\varepsilon = 0.10$.

In an early implementation, an additive constant offset was used as an
``exploration term''; this is mathematically inert (it adds equally to
all neighbors and does not change the argmin), and was identified by
an adversarial audit. The $\varepsilon$-greedy formulation replaces
it with a real exploration mechanism.

\subsection{Implicit Momentum}

The visited set implements infinite inertia by construction --- a
previously visited node is never reconsidered, equivalent to assigning
it infinite potential. Empirical tests with explicit momentum terms
confirmed no measurable additional effect, so explicit momentum is
omitted.

\section{Experimental Setup}

\subsection{Topologies}

\begin{itemize}[itemsep=1pt]
  \item \textbf{Synthetic}: Erd\H{o}s--R\'enyi $G(N, p)$ with
        size-dependent density coverage:
    \begin{itemize}[itemsep=1pt]
      \item $N = 20$: $p \in \{0.05, 0.10, 0.15, 0.20, 0.25, 0.30, 0.40, 0.50\}$
      \item $N = 50$: $p \in \{0.05, 0.10, 0.15, 0.20, 0.25, 0.30\}$
      \item $N = 100$: $p \in \{0.05, 0.10, 0.15, 0.20\}$ (sparse
            regime only, where the phase transition is observable;
            denser $n{=}100$ graphs converge trivially to zero loss
            for all strategies)
    \end{itemize}
  \item \textbf{Real}: Abilene ($N = 11$, $\rho = 0.25$) and GEANT
        ($N = 40$, $\rho = 0.08$) from the Internet Topology
        Zoo~\cite{knight2011}
\end{itemize}

\subsection{Traffic Model}

For each scenario and seed, 200 random source--destination pairs are
generated from a fixed traffic seed, independent of any routing
strategy's randomness. All strategies see identical traffic.

\subsection{Statistical Validation}

100 independent seeds per scenario for synthetic $n{=}20$, $n{=}50$,
and real topologies; 50 seeds for $n{=}100$ due to computational cost.
Mean and standard deviation reported throughout. Across the three
batteries (canonical 4-strategy, lookahead $h{=}2$ with 6 strategies,
and adaptive-$\beta$ with 6 strategies), 28{,}800 simulations were
executed ($1800 \times 4 + 1800 \times 6 + 1800 \times 6$).

\subsection{Baselines}

\begin{itemize}[itemsep=1pt]
  \item \textbf{SP}: Dijkstra on $\mathrm{latency}$
  \item \textbf{ECMP}: random choice among shortest
        paths~\cite{rfc2992}
  \item \textbf{LASP}: Dijkstra on
        $\mathrm{latency} \cdot (1 + \mathrm{load}/\mathrm{capacity})$
\end{itemize}

\subsection{Audit-Clean Implementation}

Three independent adversarial code reviews identified and corrected:
\begin{itemize}[itemsep=1pt]
  \item RNG contamination across strategies (separate RNGs for
        traffic, baseline path choice, and EMMET exploration)
  \item Asymmetric information in EMMET's loss term (snapshot is
        read-only during measurement; populated only by warm-up)
  \item Decay scale mismatch (switched from $0.95^{\mathrm{step}}$ to
        half-life formulation)
  \item Mathematically inert exploration constant (replaced with
        $\varepsilon$-greedy)
  \item Non-portable absolute file paths (replaced with portable
        \texttt{Path(\_\_file\_\_).parents[1]})
  \item Latency metric mixing partial work of lost packets (separated
        into latency-per-delivered and latency-per-attempt)
\end{itemize}
All reported results are from the audited implementation.

\subsection{Selection-Bias Verification}\label{sec:bias}

The metric \emph{latency per delivered packet} is well-defined within
each strategy but, when compared across strategies that deliver
different subsets of packets, is exposed to selection bias: a strategy
that drops more packets may show artificially low mean latency because
the surviving packets are not a random sample. To rule out this
artefact, we re-ran a focused experiment on $\mathrm{ER}(20, p)$ for
$p \in \{0.05, 0.10, 0.15, 0.20, 0.25, 0.30\}$ with 50 seeds per
density, recording per-packet outcomes for SP, LASP and EMMET full. We
then compared two metrics: (i) the original mean latency over each
strategy's own delivered set, and (ii) the \emph{conditional} mean
latency over the subset of packets delivered by all three strategies
(the common delivery set). The two metrics differ by less than 0.9
latency units across all six densities, and the strategy ordering
(SP fastest, LASP intermediate, EMMET slowest) is identical under both
metrics. Full per-density numbers are in
\texttt{data/selection\_bias\_analysis.json}. We therefore retain
\emph{latency per delivered packet} as the headline latency metric.

\subsection{Default Parameters}

$\alpha = 1.0$, $\beta_{\mathrm{base}} = 3.0$, $\gamma = 2.0$,
$\varepsilon = 0.10$, $\theta = 1.0$, $\mathrm{TTL\_FACTOR} = 2$,
$\mathrm{HALF\_LIFE} = 100$ steps.

\section{Results}

\subsection{Density Sweep: Phase Transition}

Density swept on Erd\H{o}s--R\'enyi $G(20, p)$ from $p{=}0.05$ to
$p{=}0.50$ (Table~\ref{tab:density}).

\begin{table}[h]
\centering
\caption{Density sweep on $G(20, p)$, 100 seeds per point.}
\label{tab:density}
\begin{tabular}{lrrrr}
\toprule
Density & SP loss & LASP loss & EMMET full loss & Connectivity \\
\midrule
0.05 & 11.98 & 11.53 & \textbf{5.15}  & 0\,\% \\
0.10 & 39.67 & 35.14 & \textbf{12.30} & 7\,\% \\
0.15 & 18.81 & 11.48 & \textbf{7.61}  & 43\,\% \\
0.20 & 10.57 & 4.78  & \textbf{2.39}  & 80\,\% \\
0.25 & 4.57  & 1.14  & 0.41           & 93\,\% \\
0.30 & 2.45  & 0.53  & 0.17           & 100\,\% \\
0.40 & 1.07  & 0.18  & 0.04           & 100\,\% \\
0.50 & 0.83  & 0.11  & 0.00           & 100\,\% \\
\bottomrule
\end{tabular}
\end{table}

\textbf{Finding 1 --- Phase transition at $\rho_c \in [0.15, 0.30]$.}
Below this range, graphs are rarely connected; above $\rho = 0.30$,
graphs are 100\,\% connected and EMMET achieves a mean loss below 0.2 packets per 200-step run.

\textbf{Finding 2 --- Loss reduction at low density, with a latency
cost.} At $p < 0.15$, EMMET reduces losses by 54--67\,\% versus SP.
SP retains a latency advantage of 3--13\,\% (lower mean latency per
delivered packet, verified under both per-strategy and
conditional-on-common-delivery metrics; see \S\ref{sec:bias} and
\texttt{data/selection\_bias\_analysis.json}). EMMET and SP thus
occupy distinct corners of the loss-versus-latency trade-off space;
they are not strictly comparable.

\textbf{Finding 3 --- Dead ends as field collapse signature.}
Dead-end counts drop sharply at $p \approx 0.25$--$0.30$, coinciding
with full topological connectivity.

\subsection{Beta Sensitivity (Fixed \texorpdfstring{$\beta$}{beta})}

Beta swept on $G(20, 0.30)$ from $\beta{=}0.1$ to $\beta{=}5.0$ with
$\gamma{=}2.0$ (Table~\ref{tab:beta}).

\begin{table}[h]
\centering
\caption{Beta sensitivity on $G(20, 0.30)$.}
\label{tab:beta}
\begin{tabular}{lrr}
\toprule
$\beta$ & EMMET losses & EMMET latency \\
\midrule
1.0 & 0.17          & 4.856 \\
2.0 & 0.07          & 4.870 \\
3.0 & 0.07          & 4.891 \\
3.5 & \textbf{0.00} & 4.898 \\
4.0 & \textbf{0.00} & 4.915 \\
4.5 & 0.03          & 4.936 \\
5.0 & 0.03          & 4.960 \\
\bottomrule
\end{tabular}
\end{table}

\textbf{Finding 4 --- $\beta$ sweet spot at 3.5--4.0.} Zero packet loss
at minimum latency cost.

\textbf{Finding 5 --- Field saturation above $\beta = 4.0$.} Excessive
congestion aversion forces packets onto longer routes that congest
previously uncongested links --- the $\beta$-loss relationship is
non-monotonic above the sweet spot. This motivates the adaptive
thermostat (\S\ref{sec:thermostat}) which sets
$\beta_{\mathrm{eff}}$ dynamically based on global load.

\subsection{Multi-Mechanism Comparison}

Six strategies evaluated on representative scenarios
(Table~\ref{tab:multi}).

\begin{table}[h]
\centering
\caption{Multi-mechanism comparison: mean packet loss, 100 seeds (50
for $n{=}100$). Lower is better. ``cold'' denotes greedy potential field
without snapshot; ``thermal'' adds warm-up + decay +
$\varepsilon$-greedy; ``adaptive'' uses adaptive-$\beta$ without
snapshot; ``full'' combines adaptive-$\beta$ with thermal.}
\label{tab:multi}
\small
\begin{tabular}{lrrrrrr}
\toprule
Scenario & SP & LASP & cold & thermal & adaptive & full \\
\midrule
ER $\rho{=}0.05$ ($n{=}20$) & 11.98 & 11.53 & 11.25 & 5.46    & 10.96    & \textbf{5.15}  \\
ER $\rho{=}0.10$ ($n{=}20$) & 39.67 & 35.14 & 33.67 & 13.15   & 32.24    & \textbf{12.30} \\
ER $\rho{=}0.05$ ($n{=}50$) & 16.87 & 11.70 & 12.82 & 5.27    & 12.10    & \textbf{4.50}  \\
Abilene                     & 56.67 & 46.37 & 44.36 & 43.68   & 42.62    & \textbf{40.65} \\
GEANT                       & 26.38 & 11.68 & 11.11 & 6.05    & 9.97     & \textbf{5.31}  \\
\bottomrule
\end{tabular}
\end{table}

\textbf{Finding 6 --- Thermal memory is the dominant single mechanism.}
Adding the warm-up snapshot with decay reduces losses by 46\,\% to
59\,\% over greedy-cold across the board.

\textbf{Finding 7 --- Adaptive $\beta$ contributes additional gains
synergistically with thermal memory.} EMMET full (both mechanisms)
outperforms either alone on every scenario tested. The full system
reduces losses by 55\,\% to 65\,\% versus LASP on synthetic networks
and 54.5\,\% on GEANT.

\textbf{Finding 8 --- Adaptive $\beta$ is the only mechanism that
improves Abilene.} The small Abilene topology ($N{=}11$) was the
difficult case for previous variants. Adaptive $\beta$ alone reduces
losses by 8.1\,\% versus LASP, and the full system reaches 12.3\,\%
--- the first significant improvement on this scenario.

\textbf{Finding 9 --- ECMP degenerates to SP without equal-cost
ties.} With real-valued latency, ECMP rarely finds multiple equal-cost
paths and produces results identical to SP. ECMP is effective only in
topologies engineered with uniform link costs (e.g., datacenter
fat-trees).

\subsection{Mechanism Interaction: Synergy vs.\ Substitution}

A separate battery
(\texttt{experiments/emmet\_lookahead.py}, results in
\texttt{data/lookahead\_summary.json}) tested local lookahead with
horizon $h{=}2$: at each routing decision, the potential is evaluated
over the next two hops rather than one (Table~\ref{tab:lookahead}).

\begin{table}[h]
\centering
\caption{Lookahead $h{=}2$ effect on EMMET cold and thermal variants.}
\label{tab:lookahead}
\begin{tabular}{lrrrr}
\toprule
Scenario & cold & cold + LA2 & thermal & thermal + LA2 \\
\midrule
ER $\rho{=}0.05$ ($n{=}20$) & 11.25 & 11.37 & 5.46  & 7.85  \\
ER $\rho{=}0.10$ ($n{=}20$) & 33.67 & 34.10 & 13.15 & 19.13 \\
ER $\rho{=}0.05$ ($n{=}50$) & 12.82 & 12.18 & 5.27  & 7.68  \\
GEANT                       & 11.11 & 10.42 & 6.05  & 6.19  \\
\bottomrule
\end{tabular}
\end{table}

\textbf{Finding 10 --- Local lookahead does not improve EMMET; it
degrades the thermal variant.} Lookahead $h{=}2$ yields negligible
change to EMMET cold (within $\pm 1.0\,\%$ in $n{=}20$ sparse and
modest improvement of 5--6\,\% in $n{=}50$ and GEANT) but consistently
increases losses for EMMET thermal by 2\,\% to 45\,\%. The combination
is therefore counterproductive whenever a thermal snapshot is in use.
We interpret this as a substitution effect: thermal memory and local
lookahead both provide non-myopic information about near-term link
state, and their signals interfere when combined. The snapshot is a
more efficient signal for the same purpose.

This contrasts with adaptive $\beta$, which is \textbf{complementary}
to thermal memory. The distinction is meaningful: synergistic
mechanisms operate on different signals ($\beta$ operates on global
state, thermal on per-edge history), while substitutive mechanisms
compete for the same signal (both lookahead and thermal aim to
anticipate failure).

\subsection{Pareto Frontier Analysis}

Pareto analysis (\texttt{experiments/emmet\_pareto.py}, results in
\texttt{data/pareto\_summary.json}). Across the 20 scenarios,
Pareto-optimal frequency in the (latency-per-delivered, mean-loss)
plane (Table~\ref{tab:pareto}).

\begin{table}[h]
\centering
\caption{Pareto-optimal frequency across 20 scenarios.}
\label{tab:pareto}
\begin{tabular}{lrr}
\toprule
Strategy        & Pareto-optimal & \% of scenarios \\
\midrule
SP              & 20 / 20        & 100\,\% \\
LASP            & 17 / 20        & 85\,\%  \\
EMMET cold      & 14 / 20        & 70\,\%  \\
EMMET thermal   & 3 / 20         & 15\,\%  \\
EMMET adaptive  & 8 / 20         & 40\,\%  \\
EMMET full      & 8 / 20         & 40\,\%  \\
\bottomrule
\end{tabular}
\end{table}

\textbf{Finding 11 --- SP and EMMET full occupy distinct corners of
the Pareto frontier.} SP is always Pareto-optimal because it minimizes
latency; EMMET full is Pareto-optimal in 40\,\% of scenarios because
it minimizes loss. The Pareto analysis confirms EMMET as the
loss-minimizing extreme of a routing trade-off, not as a universal
dominator.

\section{Discussion}

\subsection{Phase Transition Interpretation}

The transition at $\rho_c \in [0.15, 0.30]$ has a natural physical
interpretation. Below the critical density, the graph is frequently
disconnected; the potential field has no gradient to follow because
the destination is unreachable from many source nodes. This is
analogous to a system below its percolation threshold. Above the
critical density, the field is well-defined everywhere.

\subsection{Field Saturation and the Adaptive Thermostat}

The non-monotonic $\beta$ behavior above $\beta = 4.0$ reveals a
saturation effect: overly aggressive congestion aversion forces
packets onto longer paths, saturating links that would otherwise have
been free. The adaptive thermostat resolves this by setting
$\beta_{\mathrm{eff}}$ dynamically: high during stress, low during
relaxation. This converts a parameter-tuning problem into a
self-regulating mechanism.

\subsection{Synergy and Substitution Among Mechanisms}

The empirical distinction between synergistic and substitutive
mechanisms is one of the main contributions of this work. A mechanism
designer adding components to an algorithm typically assumes additive
benefit; our results show this assumption fails for some mechanism
pairs. The distinction appears to follow from whether the mechanisms
operate on overlapping or distinct signals:
\begin{itemize}[itemsep=1pt]
  \item \textbf{Adaptive $\beta$ + thermal memory} --- distinct
        signals (global load vs.\ per-edge history) $\rightarrow$
        synergy.
  \item \textbf{Lookahead $h{=}2$ + thermal memory} --- overlapping
        purpose (anticipate failure) $\rightarrow$ substitution.
\end{itemize}
This suggests a design principle: when adding a mechanism to a
potential-field algorithm, identify the signal it consumes and check
whether existing mechanisms already exploit that signal.

\subsection{Limitations}\label{sec:limits}

\begin{itemize}[itemsep=1pt]
  \item $N \leq 100$ nodes synthetic; real-Internet-scale topologies
        ($> 1000$ nodes) not yet measured
  \item Warm-up phase is sensitive to the ratio of warmup\_steps to
        $|V|$ in very small topologies
  \item Static $\alpha$, $\gamma$, $\varepsilon$; adaptive scheduling
        deferred to future work
  \item All results from simulation; no real deployment tested
  \item LASP is a strong but not exhaustive baseline; comparison
        against TE controllers (e.g.\ MATE, B4-style flow assignments)
        deferred to future work
\end{itemize}

\section{Conclusion}

We presented EMMET, a physics-inspired adaptive routing algorithm
modeling packets as particles in a composite potential field with
thermal dynamics and an adaptive global thermostat. We characterized
empirically:
\begin{enumerate}[itemsep=2pt]
  \item A phase transition at $\rho_c \in [0.15, 0.30]$ below which
        the field collapses
  \item A clear loss-versus-latency trade-off across the density range, with EMMET dominating loss and SP dominating latency
  \item A $\beta$ sweet spot at 3.5--4.0 followed by a field
        saturation regime
  \item Synergy between adaptive $\beta$ and thermal memory
  \item Substitution between local lookahead and thermal memory
  \item Loss reductions of 55--65\,\% versus LASP on synthetic
        networks and 54.5\,\% on GEANT, with 12.3\,\% reduction on
        the small Abilene topology that previous variants could not
        improve.
\end{enumerate}

EMMET trades a modest latency cost for substantial loss reduction.
The physical framework provides interpretable parameters with clear
analogues --- friction, heat, energy, thermostat, exploration --- and
identifies failure regimes (field collapse, saturation) and mechanism
interactions (synergy, substitution) that prior adaptive routing
literature has not formally characterized.

Three independent adversarial code reviews validated the
implementation. All data, code, and reproducibility scripts are
available at the project repository under MIT license.

\section*{Code and Data Availability}

All experiments are reproducible from the public repository:
\url{https://github.com/Carloscodix/EMMET}. The repository includes
implementation source, experiment scripts, raw and aggregated JSON
results for all 28{,}800 simulations, paper-grade figures, and
\texttt{requirements.txt} for environment setup. License: MIT.

\begin{thebibliography}{9}

\bibitem{lenders2008}
V.~Lenders, M.~May, B.~Plattner.
\newblock Density-based anycast: A robust routing strategy for wireless
  ad hoc networks.
\newblock \textit{IEEE/ACM Transactions on Networking}, 2008.

\bibitem{tassiulas1992}
L.~Tassiulas, A.~Ephremides.
\newblock Stability properties of constrained queueing systems and
  scheduling policies for maximum throughput.
\newblock \textit{IEEE Trans.\ Automatic Control}, 1992.

\bibitem{karp2000gpsr}
B.~Karp, H.~T.~Kung.
\newblock GPSR: Greedy perimeter stateless routing for wireless
  networks.
\newblock In \textit{Proc.\ MobiCom}, 2000.

\bibitem{rfc2992}
C.~Hopps.
\newblock Analysis of an Equal-Cost Multi-Path Algorithm.
\newblock \textit{RFC 2992}, IETF, 2000.

\bibitem{knight2011}
S.~Knight, H.~X.~Nguyen, N.~Falkner, R.~Bowden, M.~Roughan.
\newblock The Internet Topology Zoo.
\newblock \textit{IEEE Journal on Selected Areas in Communications},
  2011.

\end{thebibliography}

\end{document}
